% ============================================================================
% Fichier  : partie2/P02_C05_nist_iso.tex
% Titre    : Le NIST Privacy Framework et les Normes ISO
% Auteur   : MINKA MI NGUIDJOÏ Thierry Emmanuel
% Version  : 3.0 -- novembre 2026
% Corrigés : corriges/Corr_C05.tex
% ============================================================================
% Sources : NIST Privacy Framework 1.0 (16 janvier 2020) et projet public de
% la version 1.1 (14 avril 2025) ; NISTIR 8062 ; SP 800-53 rév. 5 ;
% SP 800-122 ; AI RMF 1.0. Normes ISO/IEC 29100, 27701 (éditions 2019 et
% 2025), 29134, 27018, sous licence, paraphrasées. Les intitulés anglais des
% fonctions du NIST sont conservés. État de la version 1.1 du cadre du NIST
% à revérifier à chaque édition du cours. Les sections sur les normes ISO et
% les deux certifications viennent de l'ancien chapitre 9 (version 1).
% ============================================================================

\chapter{Le NIST Privacy Framework et les Normes ISO}
\label{chap:normes}
\chaptermark{NIST et normes ISO}

\epigraph{%
    « Un certificat dit qu'un système de management existe.\\
    Il ne dit pas que chaque traitement est licite. »%
}{Principe du cours}

\niveauChapitre{Fondamental}{%
    Chapitres~\ref{chap:international} et~\ref{chap:rgpd}~; ISO/IEC~27001
    (UE MSN-2S3-SEC08)
}

\begin{boxobjectifs}
\begin{itemize}
    \item Distinguer trois natures de référentiel~: la loi, la norme
          certifiable, le cadre volontaire.
    \item Décrire le \textit{Privacy Framework} du NIST~: noyau, profils,
          niveaux de mise en œuvre.
    \item Situer les normes ISO/IEC~29100, 27701, 29134 et 27018 et dire ce
          que l'édition 2025 de l'ISO/IEC~27701 a changé.
    \item Distinguer la certification prévue par la loi (article~34) de la
          certification d'un système de management.
    \item Établir une correspondance entre la loi, le RGPD, le cadre du NIST
          et la norme et en dire les limites.
\end{itemize}
\end{boxobjectifs}

\bigskip

% ============================================================================
\section*{Avant de commencer : « avec le certificat, nous serons conformes »}
% ============================================================================

Bafoussam, mardi 1\textsuperscript{er}~septembre~2026. NOUN PAY\footnote{NOUN
PAY, ses dirigeants et ses clients sont fictifs. L'entreprise figure aussi
dans le cours de gouvernance des SI de l'auteur.} prépare une levée de fonds.
Un investisseur américain demande une preuve de conformité en matière de
données et joint à son questionnaire une grille tirée du \textit{Privacy
Framework} du NIST. Le directeur général, M.~NJOYA, a trouvé la solution~: « Nous
allons nous faire certifier ISO/IEC~27701. Avec le certificat, nous serons conformes. »

Mme~BELLO, sollicitée pour un diagnostic, lui pose deux questions. Le
périmètre du certificat couvrira-t-il la notation automatique des
crédits, ou seulement l'application de paiement~? Et l'organisme
certificateur vérifiera-t-il que chaque traitement a reçu l'autorisation de
l'article~19~? À la seconde, la réponse est non~: un auditeur de
certification vérifie le système de management contre la norme, il ne
rend pas un avis juridique. Quant à la grille du NIST, personne à NOUN PAY
ne sait ce qu'elle mesure.

\begin{boxalert}
\textbf{Violation typique.} Présenter un certificat de système de
management comme une preuve de conformité à la loi.

\smallskip
\textbf{Règle~:} un certificat ISO/IEC~27701 atteste qu'un système de
management de la protection des données existe et fonctionne sur un
périmètre donné. La conformité à la loi se démontre par l'audit de
conformité légale, exigence par exigence. Les deux se complètent~; l'un ne
remplace pas l'autre.
\end{boxalert}

% ============================================================================
\section{Trois Natures de Référentiel}
\label{sec:c05-natures}
% ============================================================================

Les textes rencontrés depuis le début de cette partie n'ont pas la même
nature. L'auditeur les range en trois familles avant de s'en servir.

\begin{tableaucours}{L{2.6cm}L{3.2cm}L{3.0cm}L{3.2cm}}{%
    \tableauHeader{Nature} &
    \tableauHeader{Exemples} &
    \tableauHeader{Qui l'impose} &
    \tableauHeader{Ce que prouve la conformité}
}
Loi & Loi n°~2024/017, loi n°~2010/012, RGPD & L'État, sous peine de
sanction & Le respect d'obligations \\
Norme certifiable & ISO/IEC~27701, ISO/IEC~27001 & Un client, un marché,
l'organisation elle-même & L'existence d'un système de management, attestée
par un tiers \\
Cadre volontaire & \textit{Privacy Framework} du NIST & Personne~; un
partenaire peut le demander & Une maturité déclarée, sans certificat \\
\end{tableaucours}
\vspace{-0.8em}
\begin{center}{\small\itshape Tableau~5.1. Trois natures de référentiel}\end{center}

Un cadre volontaire devient contraignant de la même manière qu'un texte
étranger~: par contrat. L'investisseur de NOUN PAY peut faire du profil NIST
une condition de son entrée au capital. Le cadre reste alors contractuel au
sens du Chapitre~\ref{chap:international}.

% ============================================================================
\section{Le Privacy Framework du NIST}
\label{sec:c05-nist}
% ============================================================================

Le \textit{National Institute of Standards and Technology} est l'institut de
normalisation du département américain du commerce. Il a publié en
janvier~2020 la version~1.0 de son \textit{Privacy Framework}~\cite{nist2020pf},
un outil de gestion du risque d'atteinte à la vie privée. Le cadre est
volontaire, gratuit et sans certification. Il ne dépend d'aucune loi~: il a
été conçu dans un pays qui n'a pas de loi générale de protection des données
(Chapitre~\ref{chap:international}), pour des organisations qui doivent en
respecter plusieurs à la fois.

Il se compose de trois éléments, repris du cadre de cybersécurité du même
institut.

\subsection{Le noyau}

Le \termecle{noyau} (\textit{Core}) est un catalogue de résultats attendus,
rangés en cinq fonctions. Chaque fonction se divise en catégories, puis en
sous-catégories rédigées comme des résultats et non comme des moyens.

\begin{tableaucours}{L{2.8cm}L{5.0cm}L{4.2cm}}{%
    \tableauHeader{Fonction} &
    \tableauHeader{Question posée (paraphrasée)} &
    \tableauHeader{Écho dans la loi n°~2024/017}
}
\textit{Identify-P} & Quelles données traite-t-on, pour qui, avec quel
risque pour les personnes~? & Registre (art.~29), évaluation du risque
(art.~25, 33) \\
\textit{Govern-P} & Qui décide, selon quelles règles et quelles
obligations~? & Responsabilité des acteurs (art.~20), autorisation
(art.~19) \\
\textit{Control-P} & Sait-on gérer les données finement~: accès, correction,
suppression, conservation~? & Droits (art.~37 à~46), durée (art.~13) \\
\textit{Communicate-P} & Les personnes comprennent-elles ce que l'on fait de
leurs données~? & Information (art.~14, 21) \\
\textit{Protect-P} & Les données sont-elles protégées contre les
incidents~? & Sécurité (art.~27), violations (art.~22) \\
\end{tableaucours}
\vspace{-0.8em}
\begin{center}{\small\itshape Tableau~5.5. Les cinq fonctions du noyau}\end{center}

La fonction \textit{Protect-P} recouvre le cadre de cybersécurité du NIST,
dont les fonctions de détection, de réponse et de reprise complètent le
dispositif en cas de violation. Le cadre distingue ainsi deux risques que ce
cours a séparés dès le Chapitre~\ref{chap:principes}~: celui qui vient d'un
incident de sécurité et celui qui vient du traitement lui-même, même
parfaitement sécurisé. Le NIST nomme le second « action problématique sur
les données ».

\subsection{Les profils}

Un \termecle{profil} est la sélection des résultats du noyau qui comptent
pour une organisation. On en dresse deux. Le profil actuel décrit ce qui est
atteint aujourd'hui. Le profil cible décrit ce que l'organisation veut
atteindre, compte tenu de ses obligations et de ses risques. L'écart entre
les deux donne le plan d'action.

Le profil actuel se prête à la cotation du cours. Chaque résultat retenu
devient une proposition fermée~: atteint (\OUI{}), non atteint (\NON{}), ou
pas assez d'éléments pour le dire (\NEK{}). Un profil que l'organisation
remplit seule, sans preuve, vaut une déclaration au sens du
Chapitre~\ref{chap:audit}.

\subsection{Les niveaux de mise en œuvre}

Les \termecle{niveaux} (\textit{Tiers}) décrivent la manière dont
l'organisation gère le risque, de~1 à~4~: partielle, informée par le risque,
reproductible, adaptative. Ce ne sont pas des niveaux de maturité à
atteindre tous. Une petite structure peut viser légitimement le niveau~2.

\begin{boxinfo}
\textbf{La version 1.1.} Le NIST a publié en avril~2025 le projet d'une
version~1.1, qui réaligne le cadre sur la version~2.0 de son cadre de
cybersécurité et ajoute des éléments sur l'intelligence artificielle. À la
date de ce cours, l'institut annonçait la version finale pour 2026. La
structure en noyau, profils et niveaux ne change pas. L'auditeur à qui l'on
présente un profil vérifie la version qu'il utilise.
\end{boxinfo}

\subsection{Les autres publications utiles}

\begin{tableaucours}{L{3.0cm}L{4.8cm}L{4.2cm}}{%
    \tableauHeader{Publication} &
    \tableauHeader{Objet (paraphrasé)} &
    \tableauHeader{Usage dans ce cours}
}
NISTIR 8062 & Ingénierie de la vie privée~: prévisibilité, maîtrise,
dissociabilité & Évaluer le risque (Chapitre~\ref{chap:risque}) \\
SP 800-53, révision 5 & Catalogue de mesures de sécurité et de vie privée &
Détailler une mesure (Chapitre~\ref{chap:securite}) \\
SP 800-122 & Protection de la confidentialité des données
d'identification & Classer les données par niveau d'impact \\
SP 800-61 & Gestion des incidents & Violations
(Chapitre~\ref{chap:violations}) \\
AI RMF 1.0 & Gestion des risques de l'intelligence artificielle & Profilage
(Chapitre~\ref{chap:profilage}) \\
\end{tableaucours}
\vspace{-0.8em}
\begin{center}{\small\itshape Tableau~5.6. Cinq publications du NIST pour
l'auditeur}\end{center}

Les trois objectifs du NISTIR~8062 complètent utilement la triade de la
sécurité. La \textit{prévisibilité}~: la personne et l'organisation peuvent
anticiper ce qui sera fait des données. La \textit{maîtrise}~: on peut
modifier, supprimer ou communiquer une donnée à un niveau fin. La
\textit{dissociabilité}~: on peut traiter sans rattacher les données à une
personne quand ce lien n'est pas nécessaire.

\subsection{Ce que le cadre n'apporte pas}

Le cadre du NIST ne dit pas ce qui est permis. Il ne connaît ni
l'autorisation préalable, ni les trois dérogations au consentement, ni la
liste camerounaise des données sensibles. Une organisation peut présenter un
profil cible entièrement atteint et traiter des données sans fondement.

% ============================================================================
\section{La Famille des Normes}
\label{sec:c05-famille}
% ============================================================================

\begin{tableaucours}{L{2.6cm}L{5.0cm}L{4.4cm}}{%
    \tableauHeader{Norme} &
    \tableauHeader{Objet (paraphrasé)} &
    \tableauHeader{Usage dans l'audit}
}
ISO/IEC 29100 & Cadre de la protection de la vie privée~: vocabulaire,
acteurs, onze principes & Référentiel de principes, comparaison avec la
loi \\
ISO/IEC 27701 & Système de management de la protection de la vie privée &
Audit contre une norme, certification \\
ISO/IEC 29134 & Lignes directrices pour l'évaluation d'impact & Méthode de
l'évaluation d'impact (Chapitre~\ref{chap:risque}) \\
ISO/IEC 27018 & Protection des données personnelles dans le nuage public,
pour les sous-traitants & Évaluer un hébergeur
(Chapitre~\ref{chap:transferts}) \\
ISO/IEC 27001 et 27002 & Système de management et mesures de sécurité de
l'information & Domaine sécurité (Chapitre~\ref{chap:securite}) \\
ISO 19011, ISO/IEC 27007 & Conduite des audits & Méthode
(Chapitre~\ref{chap:audit}) \\
\end{tableaucours}
\vspace{-0.8em}
\begin{center}{\small\itshape Tableau~5.2. Les normes utiles à l'audit de protection des données}\end{center}

Toutes ces normes sont publiées sous licence~: ce cours les paraphrase et
n'en reproduit aucun passage. L'auditeur qui fonde un constat sur une
exigence normative doit disposer de l'édition en vigueur et la citer avec
son année.

\subsection{Les onze principes de l'ISO/IEC 29100}

L'ISO/IEC~29100~\cite{iso29100} énonce onze principes~: le consentement et le
choix~; la légitimité et la spécification des finalités~; la limitation de
la collecte~; la minimisation des données~; la limitation de l'utilisation,
de la conservation et de la divulgation~; l'exactitude et la qualité~;
l'ouverture, la transparence et l'information~; la participation et l'accès
des personnes~; la responsabilité~; la sécurité~; et le respect des
obligations de protection de la vie privée. La comparaison avec la loi
camerounaise fait apparaître la minimisation, absente comme principe
général de la loi (Chapitre~\ref{chap:principes}).

\subsection{L'ISO/IEC 27701}

L'ISO/IEC~27701~\cite{iso27701} étend la logique de l'ISO/IEC~27001 aux
données personnelles~: contexte, leadership, planification, support,
fonctionnement, évaluation et amélioration, plus des mesures propres aux
responsables de traitement et aux sous-traitants. Publiée en 2019 comme
extension de l'ISO/IEC~27001, elle a été révisée en octobre~2025 pour
devenir une norme de système de management autonome~: la certification
ISO/IEC~27001 n'est plus un préalable. Les deux éditions coexistent pendant
la période de transition~; l'auditeur vérifie l'édition sur laquelle porte
le certificat qu'on lui présente.

Elle apporte à l'audit trois choses que la loi ne donne pas~: une méthode
de management (un cycle d'amélioration, des revues de direction, des
audits internes), des mesures détaillées propres aux données personnelles
et un vocabulaire partagé avec les partenaires étrangers. Elle n'apporte
pas les obligations propres au Cameroun~: l'autorisation préalable, les
trois dérogations au consentement, le rapport annuel de l'article~27, la
liste camerounaise des données sensibles.

% ============================================================================
\section{Deux Certifications}
\label{sec:c05-certification}
% ============================================================================

\begin{boxjuris}
\textbf{\loicam{n°~2024/017}, article~34}~\cite{loi2024017}~:
\textit{« L'Autorité de protection des données à caractère personnel met en
place un mécanisme de certification du traitement des données à caractère
personnel conforme aux principes et exigences énoncés par la présente
loi. »}

\smallskip
\textbf{Portée.} L'article~5 définit la certification comme un outil
permettant de communiquer sur le niveau de protection « au regard des
critères d'un référentiel approuvé au préalable par l'Autorité ». Il s'agit
d'une certification de conformité à la loi, distincte des certifications
ISO. L'article~59 punit le responsable qui « opère sans certification
prévue dans la présente loi »~: tant que le mécanisme n'existe pas, la
portée de cette sanction est un \NEK{} de critère.
\end{boxjuris}

\begin{tableaucours}{L{3.2cm}L{4.4cm}L{4.4cm}}{%
    \tableauHeader{} &
    \tableauHeader{Certification de l'article~34} &
    \tableauHeader{Certification ISO/IEC 27701}
}
Délivrée par & L'Autorité, ou selon son mécanisme & Un organisme de
certification accrédité \\
Référentiel & Approuvé par l'Autorité & Norme internationale \\
Objet & Traitement, processus, produit & Système de management sur un
périmètre \\
Portée juridique & Prévue par la loi (articles~34, 59) & Aucune au regard de
la loi \\
État en novembre~2026 & Mécanisme non publié & Disponible \\
\end{tableaucours}
\vspace{-0.8em}
\begin{center}{\small\itshape Tableau~5.3. Deux certifications à ne pas confondre}\end{center}

% ============================================================================
\section{Correspondances entre les Référentiels}
\label{sec:c05-correspondances}
% ============================================================================

Une organisation exposée à plusieurs référentiels établit une table de
correspondance. Elle évite de faire plusieurs fois le même travail~: le
registre de l'article~29 sert aussi le RGPD, le cadre du NIST et la norme.

\begin{tableaucours}{L{2.6cm}L{2.2cm}L{2.2cm}L{2.4cm}L{2.6cm}}{%
    \tableauHeader{Thème} &
    \tableauHeader{Loi n°~2024/017} &
    \tableauHeader{RGPD} &
    \tableauHeader{NIST} &
    \tableauHeader{ISO/IEC 27701}
}
Recensement & Art.~29 & Art.~30 & \textit{Identify-P} & Registre des
traitements \\
Fondement & Art.~9 & Art.~6 & \textit{Govern-P} & Base légale identifiée \\
Autorisation & Art.~19 & Aucun & Aucun & Aucun \\
Information & Art.~21 & Art.~13, 14 & \textit{Communicate-P} & Information
des personnes \\
Droits & Art.~37 à~47 & Art.~15 à~22 & \textit{Control-P} & Obligations
envers les personnes \\
Risque et impact & Art.~25, 33 & Art.~35 & \textit{Identify-P} & Évaluation
d'impact \\
Sécurité & Art.~27 & Art.~32 & \textit{Protect-P} & Mesures de sécurité \\
Sous-traitants & Art.~30 & Art.~28 & \textit{Govern-P} & Mesures pour les
sous-traitants \\
Transferts & Art.~32 & Art.~44 à~49 & \textit{Govern-P} & Transferts,
pays \\
Violations & Art.~22 & Art.~33, 34 & \textit{Protect-P} et cadre de
cybersécurité & Gestion des incidents \\
Rapport annuel & Art.~27(2) & Aucun & Aucun & Aucun \\
\end{tableaucours}
\vspace{-0.8em}
\begin{center}{\small\itshape Tableau~5.4. Correspondances indicatives entre
quatre référentiels (thèmes paraphrasés)}\end{center}

Deux lignes n'ont de contenu que dans la première colonne~: l'autorisation
préalable et le rapport annuel. Aucun référentiel étranger ne les porte. Une
organisation qui bâtit sa conformité sur le RGPD, le NIST ou l'ISO seuls les
oubliera.

Une correspondance n'est jamais une équivalence. Une ligne dit que le
dispositif exigé par un référentiel peut porter l'exigence légale~; elle ne
dit pas qu'il le fait. L'audit de conformité vérifie toujours l'exigence
légale elle-même. Le Chapitre~\ref{chap:loi2010} le complétera par les
règles propres à la loi n°~2010/012 et donnera la règle pour choisir ses
critères.

% ============================================================================
\section{Le Diagnostic de NOUN PAY}
\label{sec:c05-cas}
% ============================================================================

\begin{boxscenario}{Mission NOUN PAY. Le diagnostic du 8 septembre 2026}
\textbf{Constatations.} Le périmètre envisagé pour la certification couvre
l'application de paiement et exclut la notation automatique « NOUN SCORE »,
« en cours de refonte ». Aucune demande d'autorisation n'a été préparée.
L'hébergeur, établi dans l'Union européenne, présente un certificat
ISO/IEC~27018.

\textbf{Analyse} (confiance élevée). Le certificat visé laisserait hors
périmètre le traitement le plus risqué de l'entreprise, celui qui
décide des crédits. Il serait exact et trompeur s'il était présenté comme
une preuve de conformité. Le certificat de l'hébergeur porte sur
l'hébergeur~: il ne dit rien des traitements de NOUN PAY ni du transfert
lui-même (Chapitre~\ref{chap:transferts}).
\end{boxscenario}

Pour répondre à l'investisseur, Mme~BELLO dresse avec l'équipe un profil
actuel de huit résultats, deux au moins par fonction qu'elle a pu examiner.

\begin{boxlecture}{NOUN PAY, profil actuel coté, 8 septembre 2026}
\begin{tableaucours}{L{2.4cm}L{5.4cm}cL{2.8cm}}{%
    \tableauHeader{Fonction} &
    \tableauHeader{Résultat attendu (paraphrasé)} &
    \tableauHeader{Cote} &
    \tableauHeader{Preuve}
}
\textit{Identify-P} & Les traitements et les flux sont inventoriés & \NON{} &
Aucun registre \\
\textit{Identify-P} & Le risque pour les personnes est évalué & \NON{} &
Aucune évaluation \\
\textit{Govern-P} & Les obligations légales sont recensées & \NON{} & Aucune
liste \\
\textit{Govern-P} & Les rôles sont attribués & \NON{} & Aucune
désignation \\
\textit{Control-P} & Les données sont supprimées à l'échéance & \NEK{} & Non
examiné \\
\textit{Communicate-P} & Les personnes sont informées des traitements &
\NON{} & Notice de l'application \\
\textit{Protect-P} & Les données sont chiffrées au repos et en transit &
\OUI{} & Configuration de l'hébergeur \\
\textit{Protect-P} & Les accès sont nominatifs et tracés & \OUI{} & Journaux
de la console \\
\end{tableaucours}

\medskip
\textbf{Décompte.} Deux \OUI{}, cinq \NON{}, un \NEK{} sur huit.
\hfill\lecture{0{,}25}{0{,}125}{1{,}30}

\medskip
\textbf{Lecture.} Les deux résultats atteints sont dans la seule fonction de
protection. NOUN PAY a traité la vie privée comme un problème de sécurité.
Le profil montre à l'investisseur ce qu'un certificat de sécurité n'aurait
pas montré.
\end{boxlecture}

% ============================================================================
\section{Ce que les Référentiels demandent à chaque Métier}
\label{sec:c05-metiers}
% ============================================================================

\begin{boxmetier}{Délégué à la protection des données}
Bâtir le système de management sur les obligations de la loi, puis le
rapprocher des référentiels demandés par les partenaires. Tenir la table de
correspondance et signaler les exigences camerounaises qu'aucun référentiel
étranger ne porte.
\end{boxmetier}

\begin{boxmetier}{Auditeur de la protection des données}
Lire le périmètre et l'édition de tout certificat, la version de tout
profil. Ne jamais coter \OUI{} une exigence légale sur la seule foi d'un
certificat ou d'un profil déclaré.
\end{boxmetier}

\begin{boxmetier}{Responsable de traitement}
Choisir le périmètre d'une certification en fonction des risques, pas de la
facilité. Un certificat qui exclut le traitement le plus risqué expose à un
reproche de présentation trompeuse.
\end{boxmetier}

\begin{boxresume}
\begin{itemize}
    \item Une loi oblige, une norme se certifie, un cadre volontaire se
          déclare. Les deux derniers ne deviennent contraignants que par
          contrat.
    \item Le \textit{Privacy Framework} du NIST comprend un noyau de cinq
          fonctions, des profils actuel et cible et quatre niveaux de mise
          en œuvre. Il ne dit pas ce qui est permis.
    \item L'ISO/IEC~27701 est depuis 2025 une norme autonome~; l'auditeur
          vérifie l'édition du certificat.
    \item La certification de l'article~34 est une certification de
          conformité à la loi, distincte d'une certification ISO.
    \item L'autorisation préalable et le rapport annuel n'ont d'équivalent
          dans aucun référentiel étranger.
\end{itemize}
\end{boxresume}

\begin{boxinfo}
\textbf{Points de contrôle.} Vous avez compris ce chapitre si vous savez~:
\begin{itemize}
    \item nommer les cinq fonctions du noyau du NIST et la question que pose
          chacune~;
    \item dire ce qu'un certificat ISO/IEC~27701 prouve et ce qu'il ne
          prouve pas~;
    \item citer deux obligations camerounaises absentes des référentiels
          étrangers.
\end{itemize}
\end{boxinfo}

\begin{boxexo}[title={\ding{45}\quad Exercice~5.1. Quelle norme \niveaubase}]
Quelle norme consultez-vous pour~: (a) conduire une évaluation d'impact~;
(b) évaluer un hébergeur en nuage~; (c) mettre en place un système de
management de la protection des données~; (d) définir le vocabulaire d'un
contrat~?
\end{boxexo}

\begin{boxexo}[title={\ding{45}\quad Exercice~5.2. Le périmètre de NOUN PAY \niveaumoyen}]
Rédigez l'avis de Mme~BELLO à M.~NJOYA sur le périmètre de la
certification, en séparant constatations et analyse. Proposez un périmètre
et un calendrier qui tiennent compte de la notation automatique.
\end{boxexo}

\begin{boxexo}[title={\ding{45}\quad Exercice~5.3. Compléter la correspondance \niveaumoyen}]
Ajoutez au Tableau~5.4 les lignes correspondant aux articles~22
(violations), 30 (sous-traitance) et 32 (transferts), avec le thème de la
norme et la couverture.
\end{boxexo}

\begin{boxexo}[title={\ding{45}\quad Exercice~5.4. Le profil cible \niveauexpert}]
L'investisseur demande à NOUN PAY un profil cible à douze mois.
\begin{enumerate}[(a)]
    \item Pour chacun des huit résultats du profil actuel, proposez la cible
          et la preuve qui permettra de la coter \OUI{}.
    \item Ajoutez deux résultats que la loi n°~2024/017 impose et que le
          noyau du NIST ne contient pas. Comment les présentez-vous à un
          investisseur qui ne connaît que le cadre du NIST~?
    \item Quel niveau de mise en œuvre NOUN PAY peut-elle viser
          raisonnablement~? Justifiez.
\end{enumerate}
\end{boxexo}

\bigskip

\begin{boxinfo}
\textbf{Pour aller plus loin.}
\begin{itemize}
    \item Le \textit{Privacy Framework} du NIST~\cite{nist2020pf}, librement
          accessible et le NISTIR~8062~\cite{nistir8062}.
    \item L'ISO/IEC~27701~\cite{iso27701} et l'ISO/IEC~29100~\cite{iso29100}.
    \item L'ISO/IEC~27018~\cite{iso27018}, pour les sous-traitants en nuage.
\end{itemize}
\end{boxinfo}

\begin{boxinfo}
\textbf{Transition.} Les sources internationales, le règlement européen et
les référentiels volontaires forment le paysage. Il reste à lire de près le
texte qui oblige toute organisation camerounaise et qui a internalisé une
grande partie de ce paysage~: la loi n°~2024/017, avec son Autorité et ses
sanctions. C'est l'objet du Chapitre~\ref{chap:loi}.
\end{boxinfo}
